\section{Limitations}\label{sec:limitations}
The feasibility result has a deliberately narrow statistical scope. Proposition~3 applies to deterministic, uniformly valid, distribution-free UCBs based only on iid bounded category losses. It does not rule out nonvacuous certification from fewer categories under pre-specified parametric or hierarchical models, informative side data, or randomized procedures. Such approaches exchange distribution-free generality for additional assumptions and require their own coverage analyses. The counts 14, 29, and 59 are optimistic necessary limits without multiplicity, not sufficient sample sizes for CRESS\@. The frozen family allocation and declared Hoeffding rule require more. Proposition~1 controls its declared candidate family within one fixed target/seed/source-view/condition cell; the study does not claim one simultaneous 95\% event over all reported cells.

The category interpretation also depends on the unit definition. The strict audit contains only three or four certification categories, and repeated images, corruption views, or support seeds do not constitute additional independent category draws. Image-unit bounds under the selected-mixture iid model therefore address a different estimand rather than a partial certificate for a new category. The benchmark archive is stratified by category, and alarms within a category share a support-fitted scorer; consequently, iid pooled alarm indicators are an idealized sensitivity assumption, not a design-based property of the archive. CRESS certifies source-domain risk under the assumptions of Proposition~1. The benchmark categories are a fixed heterogeneous collection, so their interpretation as iid draws from a source-category meta-population is a modeling assumption rather than an empirically testable fact in these data. A category unit includes its support construction as well as its held-out views. Moreover, the observed $L_c$ values are finite archive averages; interpreting their meta-population mean as the alarm probability of a fresh normal view requires the within-category archive to represent the declared deployment distribution. Applying the selected threshold to a particular target additionally requires source-risk dominance. An independent target draw from the same category meta-population supports only risk marginalized over that draw; it does not control every realized category. Support-only normalization and source-view selection establish neither condition. Matched source selection also assumes reliable condition metadata.

The empirical evidence is limited to MVTec, VisA, and MPDD\@. The controlled shift study covers Gaussian noise, blur, brightness/contrast, and JPEG compression, not the full range of temporal drift, sensor replacement, process changes, correlated failures, or deployment prevalence. The empirical FAR of 0.341 is the largest dataset-level aggregate in this frozen synthetic-corruption grid, not a population failure probability or a universal response to Gaussian noise. CRESS also assumes access to a separately declared archive of known-normal, held-out views from non-target source categories; in these benchmarks those views come from the normal evaluation partition. This source-assisted setup is therefore not the conventional target-category training-only benchmark protocol, although target test labels and images remain excluded from threshold construction. The image-unit sensitivity additionally assumes independent source-image draws; repeated or correlated production images would reduce its effective evidence. Prospective factory evidence is needed to determine whether the same failure mode and category-count assumptions hold under real operational shifts.

Finally, the frozen DINOv2 PCA residual is an inherited ranking substrate. Reported AnomalyDINO-S and WinCLIP results follow different protocols and serve as external references; the paper does not claim ranking state of the art. LOIO calibration and full-support test residuals remain asymmetric, and the pooled discrete-grid audit diagnoses rather than restores exchangeability. We therefore restrict the efficiency comparison to explicitly accounted category-specific storage and make no runtime claim.
